\section{Truncation exponents prime to the residue characteristic}
\label{sec:prime-to-p-all-dvr}

The divisor-class obstruction extends from dual numbers to every
truncation whose exponent is invertible in the DVR.  We use the
generic-hypersurface equivalence to define the detector on all
coefficient-free lattices.  The order of the detected class is proved
by explicit exact sequences.

\begin{theorem}\label{thm:prime-to-p-all-dvr}
Let \(O\) be a characteristic-zero DVR with residue characteristic
\(p>0\), and let \(n\ge2\) satisfy \(p\nmid n\).  Put
\(A_n=O[t]/(t^n)\).  There is a class of exact order \(n\) in
\[
 \ker\bigl(K_0(\Lat_{A_n}(S_p))
                   \longrightarrow K_0(\Lat_O(S_p))\bigr)
\]
which is zero after extension to \(\operatorname{Frac}(O)\).
An additive detector defined on the whole Grothendieck group
annihilates all Young permutation lattices and detects this class.
Thus \(\Phi_{A_n}\) is not surjective in degree \(p\).
There is no restriction on the ramification or the residue field,
and \(O\) need not be complete.
\end{theorem}

The theorem includes \(p=2\) with \(n\) odd.  It asserts the exact
order of a cyclic subgroup, without asserting that this subgroup
is a direct summand.

\subsection{Finite lattices and the mesh exact sequences}

Fix a generator \(\sigma\) of \(P=C_p\), write \(X=\sigma-1\), and set
\[
 g(X)=(1+X)^p-1=Xh(X),\qquad
 h(X)=p+\binom p2X+\cdots+X^{p-1}.
\]
For \(0\le i\le n\), define \(M_i\) over \(A_n[X]\) by generators
\(a,b\) and relations
\begin{equation}\label{eq:prime-to-p-presentations}
 Xa=t^{n-i}b,\qquad t^ia+h(X)b=0,
 \qquad
 \Phi_i=
 \begin{pmatrix}X&t^i\\-t^{n-i}&h(X)\end{pmatrix}.
\end{equation}
The determinant of \(\Phi_i\) is \(g(X)+t^n\), which is \(g(X)\)
over \(A_n\).  Its adjugate shows that \(g\) annihilates \(M_i\).

\begin{lemma}\label{lem:prime-to-p-mesh}
Each \(M_i\) is a \(P\)-lattice free of rank \(p\) over \(A_n\).
One has \(M_0\simeq M_n\simeq A_n[P]\), and, for \(1\le i<n\),
there is a coefficient-split exact sequence
\begin{equation}\label{eq:prime-to-p-mesh}
 0\longrightarrow M_i
 \xrightarrow{(\beta_{i-1},\alpha_i)}
 M_{i-1}\oplus M_{i+1}
 \xrightarrow{(\alpha_{i-1},-\beta_i)}
 M_i\longrightarrow0.
\end{equation}
Consequently, with \(x_i=[M_i]-[A_n[P]]\),
\begin{equation}\label{eq:prime-to-p-mesh-order}
 x_i=i x_1,\qquad n x_1=0,\qquad x_{n-1}=-x_1.
\end{equation}
For \(1\le i<n\), these classes vanish after coefficient reduction
\(t\mapsto0\) and after extension to \(\operatorname{Frac}(O)\).
\end{lemma}

\begin{proof}
An \(A_n\)-basis is
\begin{equation}\label{eq:prime-to-p-basis}
 a,\qquad X^j b\quad(0\le j\le p-2).
\end{equation}
The first relation specifies \(Xa\), while the second specifies
\(X^{p-1}b\), since \(h\) is monic.  Conversely these prescriptions
define an \(X\)-operator on the free module with basis
\eqref{eq:prime-to-p-basis}; both displayed relations hold there.
The resulting maps between this free module and the presented module
are inverse.  This proves freeness, and \(g(X)=0\) gives the action
of \(P\).  At \(i=0\) eliminate \(a=-hb\); at \(i=n\) eliminate
\(b=Xa\).  In both cases the remaining relation is \(g=0\), proving
the endpoint identifications.

In the bases \eqref{eq:prime-to-p-basis}, put
\[
 \alpha_i:M_i\longrightarrow M_{i+1},\quad
 \alpha_i=\operatorname{diag}(t,I_{p-1}),\qquad
 \beta_i:M_{i+1}\longrightarrow M_i,\quad
 \beta_i=\operatorname{diag}(1,tI_{p-1}).
\]
For example, \(\alpha_i\) sends \(a\) to \(ta\) and \(b\) to \(b\)
in the target, and \(\beta_i\) sends \(a\) to \(a\) and \(b\) to
\(tb\).  Substitution in \eqref{eq:prime-to-p-presentations}
proves equivariance; both composites are multiplication by \(t\).
On the first basis coordinate, the maps in
\eqref{eq:prime-to-p-mesh} are \((1,t)\) and \((t,-1)\);
on each remaining coordinate, they are \((t,1)\) and \((1,-t)\).
These formulas show that their composite is zero and that the kernel
of the second is the image of the first.  For coordinates \((r,v)\)
on the last \(M_i\), a coefficient right inverse is
\[
 (r,v)\longmapsto ((0,v),(-r,0)).
\]
A coefficient left inverse on the first map is
\[
 ((r_L,v_L),(r_R,v_R))\longmapsto(r_L,v_R).
\]
Thus \eqref{eq:prime-to-p-mesh} is admissible.  Its relations are
\(2x_i=x_{i-1}+x_{i+1}\), with \(x_0=x_n=0\), which give
\eqref{eq:prime-to-p-mesh-order}.

For \(0<i<n\), reduction modulo \(t\) is
\[
 O[X]/(X)\oplus O[X]/(h).
\]
The exact sequence
\begin{equation}\label{eq:prime-to-p-reduction}
 0\longrightarrow O[X]/(h)
  \xrightarrow{\ \cdot X\ } O[X]/(Xh)
  \xrightarrow{\ X\mapsto0\ }O\longrightarrow0
\end{equation}
is split over \(O\) by the constant polynomials on the right.
Division in \(O[X]\) verifies injectivity and the stated kernel.
This proves reduction vanishing.

Let \(K=\operatorname{Frac}(O)\).  Since \(h(0)=p\) is a unit in
\(K\), the same reduction is isomorphic to \(K[P]\).
Lift a cyclic generator of that reduction to \(M_i\otimes_OK\).
It defines a homomorphism
\[
 (K[t]/t^n)[P]\longrightarrow M_i\otimes_OK.
\]
Both sides are free of coefficient rank \(p\), and modulo \(t\)
the map is an isomorphism.  Its determinant is a unit in the local
Artinian coefficient ring, so it is an isomorphism.  This proves
generic vanishing.  None of these arguments uses completeness.
\end{proof}

\subsection{Divisor classes of exact prescribed order}

Assume for now that \(O\) is complete, with uniformizer \(\pi\)
and residue field \(k\), and set
\begin{equation}\label{eq:prime-to-p-surface}
 S=O[[X,t]]/(g(X)+t^n),\qquad B=O[[X]],\qquad
 I_a=(X,t^a)\subset S\quad(1\le a<n).
\end{equation}

\begin{lemma}\label{lem:prime-to-p-surface}
The ring \(S\) is a normal two-dimensional complete local domain.
Each \(I_a\) is rank-one reflexive and nonprincipal.  If
\(\mathfrak p=(X,t)\), then
\[
 \operatorname{div}(X)=n\mathfrak p,\qquad
 [I_a]=a[I_1]\quad\text{in }\operatorname{Cl}(S).
\]
In particular \([I_1]\) has exact order \(n\) in
\(\operatorname{Cl}(S)\).  Its class in
\(K_0(D_{\mathrm{sg}}(S)^{\mathrm{ic}})\) has order at least \(n\):
if \(r[I_1]=0\) there, then \(n\mid r\).
\end{lemma}

\begin{proof}
The polynomial \(t^n+g(X)\) is Eisenstein at \((X)\) over
\(B_{(X)}\), because \(g=Xh\) and \(h(0)=p\ne0\).
It is therefore irreducible over \(\operatorname{Frac}(B)\).
Since \(S\) is free of rank \(n\) over \(B\), it is a domain.

Its special fiber is
\[
 k[[X,t]]/(t^n+X^p),
\]
which is also a domain.  Indeed, every root of \(t^n+X^p\) over
an algebraic closure of \(k((X))\) has \(X\)-valuation \(p/n\).
A monic factor of degree \(r\) over \(k((X))\) would have constant
coefficient of valuation \(rp/n\), an integer.  Since \((n,p)=1\),
one has \(n\mid r\), proving irreducibility.  The monic
power-series quotient embeds in this field extension.

The hypersurface \(S\) is Cohen--Macaulay of dimension two.
Its only height-one prime above \(\pi\) is \((\pi)\), which has a
regular localization because its maximal ideal is principal.
At every prime not containing \(\pi\), the derivative
\[
 g'(X)=p(1+X)^{p-1}
\]
is a unit.  The derivation in \(X\) shows that the defining equation
is outside the square of the maximal ideal in the corresponding
regular ambient localization; its quotient is regular.
Thus \(S\) satisfies \(R_1\) and \(S_2\), hence is normal.

In the \(B\)-basis \(1,t,\ldots,t^{n-1}\) of \(S\), one has
\begin{equation}\label{eq:prime-to-p-ideal-basis}
 I_a=BX\oplus BXt\oplus\cdots\oplus BXt^{a-1}
            \oplus Bt^a\oplus\cdots\oplus Bt^{n-1}.
\end{equation}
This is stable under multiplication by \(t\), using \(t^n=-Xh\),
and is exactly the ideal generated by \(X,t^a\).
It is \(B\)-free and torsion free over the normal surface \(S\);
the regular sequence \(\pi,X\) proves depth two.  Thus it is
reflexive.

Its two displayed ideal generators are minimal.  If
\(cX+dt^a\in\mathfrak m_S I_a\), reduction modulo \((\pi,X)\)
in \(k[[t]]/(t^n)\) shows that \(d\in\mathfrak m_S\), because
\(t^a\notin(t^{a+1})\).  Then reduce \(cX\) modulo
\((\pi,t,X^2)\).  This quotient is \(k[X]/(X^2)\), since \(p\ge2\);
the image of \(\mathfrak m_S I_a\) is zero and that of \(X\) is
nonzero.  Hence \(c\in\mathfrak m_S\) as well.  Therefore
\(\dim_k I_a/\mathfrak m_S I_a=2\), and \(I_a\) is not principal.
This argument does not require a coefficient field in \(O\).

The only height-one prime containing \(X\) is
\(\mathfrak p=(X,t)\), with \(S/\mathfrak p=O\).  At that prime
\(h\) is a unit, \(t\) is a uniformizer, and
\(X=-t^n/h\) has valuation \(n\).  Thus
\(\operatorname{div}(X)=n\mathfrak p\).
The ideal \(I_a\) has valuation \(a\) at \(\mathfrak p\) and zero
at all other height-one primes, giving
\([I_a]=a[I_1]\), with the same ideal-class sign convention.
Nonprincipality for every \(1\le a<n\) proves exact order \(n\).
For the normality and reflexivity criteria and the identification of
rank-one reflexive modules with divisor classes, see
\citet[Tags 031S, 0AVB, and 0EBM]{Stacks2026}.

For the last assertion, the first divisor class gives a homomorphism
\[
 c_1:G_0(S)\longrightarrow\operatorname{Cl}(S).
\]
It can be defined by the determinant lattice of the torsion-free
quotient of a module, together with its torsion lengths at height-one
primes.  The determinant convention has the opposite sign to the
valuation of an ideal, while the torsion lengths have positive sign.
Localizing an exact sequence at each height-one DVR proves additivity,
and a change of generic basis changes the divisor by a principal one.
This map takes a rank-one reflexive ideal to its class and takes a
free module to zero.  Since \(S\) is local, it kills every perfect
complex and descends to \(K_0D_{\mathrm{sg}}(S)\).
The inclusion into the idempotent completion induces an injection
on \(K_0\), by \citet[Corollary 2.3, p.~6]{Thomason1997}.
Thus a relation \(r[I_1]=0\) after idempotent completion already
holds before it, and applying \(c_1\) proves \(n\mid r\).
No map from the new idempotent objects to \(\operatorname{Cl}(S)\)
is required.
\end{proof}

\subsection{Specialization of a supported formal disc}

Write
\(D_{\mathrm{sg}}(C)=D^b(\operatorname{coh}C)/\operatorname{Perf}(C)\)
and use a superscript \(\mathrm{ic}\) for idempotent completion.

\begin{lemma}\label{lem:prime-to-p-formal-disc}
Let \(S_0\) be a complete Noetherian local hypersurface, put
\(B_0=S_0[[w]]\), and let \(\mathcal T_w\) be the full subcategory
of \(D_{\mathrm{sg}}(B_0)\) consisting of objects perfect after
inverting \(w\).  There is a triangle functor
\[
 \mathcal E:\mathcal T_w^{\mathrm{ic}}
                 \longrightarrow D_{\mathrm{sg}}(S_0)^{\mathrm{ic}}
\]
which takes an \(S_0\)-module \(N\), regarded as a \(B_0\)-module
killed by \(w\), to the class of that same module \(N\).
\end{lemma}

\begin{proof}
Consider the supported quotient
\[
 \mathcal C_w=
 D^b_{V(w)}(\operatorname{coh}B_0)/
                  \operatorname{Perf}_{V(w)}(B_0).
\]
Its natural functor to \(D_{\mathrm{sg}}(B_0)\) is fully faithful.
Here is the required factorization check.  A supported bounded coherent
complex has a bounded model of coherent \(w\)-power-torsion modules;
one power \(w^N\) kills every term.  Any morphism to it from a bounded
finite-projective complex \(P^\bullet\) is represented by a chain map
and factors through \(P^\bullet/w^NP^\bullet\).  The latter is
represented by tensoring \(P^\bullet\) with the two-term Koszul
complex on \(w^N\), since \(w\) is regular on \(B_0\).  It is perfect
and supported on \(V(w)\).  The Verdier-quotient factorization
criterion now gives full faithfulness.  The supported derived model
and this criterion are \citet[Lemmas 2.1, 2.5, and 2.6,
preprint pp.~3--5]{Orlov2011};
the argument just given is affine and requires no coefficient field.

The image is dense in \(\mathcal T_w\).  Over the local hypersurface
\(B_0\), represent an object of \(\mathcal T_w\), up to shift, by a
finite maximal Cohen--Macaulay module \(L\).  The usual
matrix-factorization description of the singularity category
identifies its endomorphisms with stable module maps
\citep[equation (1.1), p.~2, and Appendix B, pp.~37--40]{BurkeWalker2012}.
Its stable endomorphism module is finite over \(B_0\): it is the
quotient of \(\operatorname{Hom}_{B_0}(L,L)\) by the image of
\(\operatorname{Hom}_{B_0}(L,B_0)\otimes_{B_0}L\).
This description commutes with localization.  After inverting \(w\),
the object is perfect, so its stable identity is zero.  Therefore
some \(w^N\) kills that identity already over \(B_0\).
The cone of \(w^N:L\to L\) has \(w\)-power-torsion cohomology,
and in the singularity category it is isomorphic to \(L\oplus L[1]\).
This proves density, hence
\[
 \mathcal C_w^{\mathrm{ic}}\simeq\mathcal T_w^{\mathrm{ic}}.
\]

Under restriction of scalars, every coherent \(w\)-power-torsion
module is finite over \(S_0\).  This gives a functor from the supported
bounded coherent category to \(D^b(\operatorname{coh}S_0)\).
A supported perfect \(B_0\)-complex has finite flat amplitude over
\(S_0\), because \(B_0\) is \(S_0\)-flat.  Its restricted cohomology
is bounded and finite, so over the Noetherian ring \(S_0\) it is
perfect.  Restriction therefore induces a triangle functor
\(\mathcal C_w\to D_{\mathrm{sg}}(S_0)\).
Extend it to idempotent completions and use the displayed equivalence.
On a module killed by \(w\) this functor is exactly the asserted
restriction of scalars.
\end{proof}

\subsection{A detector defined for every lattice}

Continue with \(O\) complete and use the surface \(S\) in
\eqref{eq:prime-to-p-surface}.  Set
\[
 Q=O[[t,X]],\qquad f=t^n,\qquad
 R=Q/(f,g)=A_n[P].
\]
Here \(f,g\) is a regular sequence.  Indeed \(t^n\) is regular,
and in \(A_n[[X]]\) multiplication by \(g\) is injective: its lowest
nonzero \(X\)-coefficient is the \(O\)-non-zero-divisor \(p\).
Distinguished-polynomial division identifies the quotient with
the finite group ring.

Form the generic hypersurface
\[
 Y=\operatorname{Proj}Q[T_f,T_g]/(fT_f+gT_g).
\]
Let \(\beta:\mathbb P^1_R\hookrightarrow Y\) and
\(\varpi:\mathbb P^1_R\to\operatorname{Spec}R\) be the natural maps.
The mixed-characteristic generic-hypersurface theorem gives mutually
inverse triangle equivalences
\begin{equation}\label{eq:prime-to-p-generic-equivalence}
 F=\beta_*\varpi^*:D_{\mathrm{sg}}(R)\longrightarrow D_{\mathrm{sg}}(Y),
 \qquad
 \Psi=R\varpi_*\beta^\flat,\qquad
 \beta^\flat=L\beta^*(-1)[-1].
\end{equation}
These are \citet[Theorem 2.10, p.~8; the paragraph after Corollary 2.11,
p.~9; and the formula preceding (4.3), p.~19]{BurkeWalker2012}.
Their hypotheses require a regular Noetherian ambient ring of finite
dimension and a regular sequence; they impose no equicharacteristic
assumption.  All those hypotheses hold here.

\begin{lemma}\label{lem:prime-to-p-detector}
There is a homomorphism
\[
 \rho:K_0(\Lat_{A_n}(P))
                  \longrightarrow K_0(D_{\mathrm{sg}}(S)^{\mathrm{ic}})
\]
defined on the full exact category.  It vanishes on \(A_n[P]\)
and on the trivial lattice \(A_n\), and satisfies
\begin{equation}\label{eq:prime-to-p-values}
 \rho([M_i])=-[I_{n-i}]\qquad(0<i<n).
\end{equation}
If \(\alpha_a\) is the automorphism \(\sigma\mapsto\sigma^a\),
\(a\in\mathbb F_p^\times\), then
\begin{equation}\label{eq:prime-to-p-normalizer-values}
 \rho(\alpha_a^*x_{n-1})=
 \begin{cases}
 -[I_1],&a=1,\\
 0,&a\ne1.
 \end{cases}
\end{equation}
\end{lemma}

\begin{proof}
We first verify that a single specialization construction is defined
for every source lattice.  If \(M\) is finite free over \(A_n\), its
twisted cyclic resolution has free \(R\)-terms \(R\otimes_{A_n}M\),
and
\[
 \operatorname{Ext}^q_R(M,M)
  =H^q(P,\operatorname{Hom}_{A_n}(M,M)),\qquad q>0.
\]
These groups are two-periodic, killed by \(p\), and of bounded finite
\(O\)-length.  The even and odd cohomology-operator modules are finite
over \(R[\chi_f,\chi_g]\), by
\citet[Remark 3.19, p.~17]{BurkeWalker2012}.  Discard finitely many
initial degrees so that \(p\) kills both modules; this does not change
their projective support.

Put \(C=R/pR\), an Artinian local ring with residue field \(k\).
The two modules are now finite graded \(C[\chi_f,\chi_g]\)-modules
whose component lengths are bounded.  Reducing modulo the maximal
ideal of \(C\) gives finite graded \(k[\chi_f,\chi_g]\)-modules with
bounded Hilbert function, hence zero-dimensional projective support.
The maximal ideal of \(C\) is nilpotent, so reduction does not change
the underlying support.  The support is therefore a finite set of
closed points of \(\mathbb P^1_k\).  This uses neither a coefficient
field in \(O\) nor any perfection assumption on \(k\).

The stable-Ext sheaf comparison identifies this support with the
nonperfect locus of \(F(M)\)
\citep[Corollaries 3.18 and 3.20, pp.~17--18;
the end of Section 8.1, pp.~33--34]{BurkeWalker2012}.
One can test nonperfectness by the stable self-End sheaf: its stalk
contains the identity precisely when the corresponding local object
is nonzero in the singularity category.  Thus \(F(M)\) has finite
nonperfect support for every lattice \(M\).

On the chart \(T_g=1\), put \(z=T_f/T_g\).  Complete at the rational
closed point \((\pi,t,X,z-1)\), and write \(w=z-1\).  The completed
ring is
\[
 \widehat Y=O[[t,X,w]]/(g(X)+(1+w)t^n).
\]
Because \(n\) is a unit, recursive division by \(n\) gives a unique
\(v(w)\in1+wO[[w]]\) with \(v(w)^n=1+w\).
The change \(T=v(w)t\) identifies this completed ring with
\[
 S[[w]],\qquad S=O[[X,T]]/(g(X)+T^n).
\]
It is the identity in the \(t\)-coordinate on \(w=0\).
The completed \(F(M)\) is perfect outside its closed point, since
its original nonperfect locus was finite and perfection is preserved
under the flat completion map.  In particular it is perfect after
inverting \(w\).  Apply Lemma~\ref{lem:prime-to-p-formal-disc}.
This defines \(\rho\).

The functor \(M\mapsto\beta_*\varpi^*M\) is exact on coherent modules:
\(\varpi\) is flat and \(\beta\) is a closed immersion.  Completion,
the coordinate change, and the functor of
Lemma~\ref{lem:prime-to-p-formal-disc} are triangle functors.
Consequently every coefficient-split exact triple of lattices gives
an additive relation under \(\rho\).  This establishes the
whole-category assertion.

\paragraph{Value on the explicit lattices.}
Over \(S\), set \(N_i=\operatorname{coker}\Phi_i\) for \(0<i<n\).
The map
\[
 a\longmapsto t^{n-i},\qquad b\longmapsto X
\]
identifies \(N_i\) with \(I_{n-i}\).
To verify the full presentation, lift a relation
\(u t^{n-i}+vX=0\) to \(O[[X,t]]\), and subtract a multiple of
\(t^n+Xh\).  The remaining relation on the regular sequence
\(t^{n-i},X\) is a multiple of \((X,-t^{n-i})\).
Thus all relations are generated by the columns
\((X,-t^{n-i})\) and \((t^i,h)\) of \(\Phi_i\).
Since \(S\) is a domain, \(f=t^n\) is regular on \(N_i\), and
the presentation gives \(N_i/fN_i=M_i\).

Let \(i_1:\operatorname{Spec}S\hookrightarrow Y\) be the section
\(T_f=T_g\), or \(z=1\).  It is globally closed and lies entirely
in the chosen chart.  On this chart \(\beta\) is cut out by \(f\).
Thus the inverse in \eqref{eq:prime-to-p-generic-equivalence} gives
\[
 \Psi(i_{1*}N_i)\simeq M_i[-1].
\]
There is no higher Tor because \(f\) is regular on \(N_i\);
\(\mathcal O(-1)\) is trivial on the section, and \(\varpi\) restricts
to an isomorphism on it.  Applying the inverse equivalence yields
\[
 F(M_i)\simeq i_{1*}I_{n-i}[1]
       \quad\text{in }D_{\mathrm{sg}}(Y).
\]
After completion this is \(I_{n-i}[1]\), killed by \(w\).
Lemma~\ref{lem:prime-to-p-formal-disc} gives
\eqref{eq:prime-to-p-values}.  The regular module \(R=A_n[P]\)
has value zero because \(\beta\) is a regular closed immersion.

The trivial \(P\)-module \(A_n=R/(X)\) also has value zero.
Near \(z=1\), \(z\) is a unit, so \(F(A_n)\) is \(Y/(X)\):
the equation modulo \(X\) forces \(zt^n=0\), hence \(f=0\).
The element \(X\) is regular on this chart.  In the regular
ambient ring, reduction of \(g+zt^n\) modulo \(X\) is the nonzero
element \(zt^n\) of the domain \(O[[t]][z]\); cancelling \(X\)
in a lifted relation proves the assertion.  Hence \(Y/(X)\)
has a length-one locally free resolution and is perfect there.

\paragraph{Normalizer separation.}
For the representative \(1\le a\le p-1\), set
\[
 \varphi_a(X)=(1+X)^a-1,\qquad
 U_a(X)=\sum_{r=0}^{a-1}(1+X)^{pr}.
\]
Then \(\varphi_a\) is an invertible formal substitution,
\(U_a\) is a unit with residue \(a\), and
\[
 g(\varphi_a(X))=U_a(X)g(X).
\]
The generic-hypersurface automorphism on the chart is
\[
 (X,z)\longmapsto(\varphi_a(X),\,U_a(X)z);
\]
it multiplies the defining equation by \(U_a\) and respects the
\(\beta,\varpi\) diagram.  Thus \(F\) is natural for this automorphism.
Pullback of support at the residual point \(z=1\) is support at
\(z=a^{-1}\).  For \(a\ne1\) that is a different rational point
over \(k\), so the completed twisted \(F(M_i)\) is perfect at the
selected point.  For \(a=1\) it is the original object.
This proves \eqref{eq:prime-to-p-normalizer-values}.
The inverse exponent results from pullback; it is not replaced
by \(a\) in this argument.
\end{proof}

\subsection{Exact order, induction, and coefficient transfer}

\begin{proof}[Proof of Theorem~\ref{thm:prime-to-p-all-dvr}]
First suppose \(O\) is complete and put \(\eta=x_{n-1}\).
Lemma~\ref{lem:prime-to-p-mesh} proves \(n\eta=0\) and both kernel
assertions.  Lemma~\ref{lem:prime-to-p-detector} gives
\(\rho(\eta)=-[I_1]\).
If \(r\eta=0\), Lemma~\ref{lem:prime-to-p-surface} therefore gives
\(n\mid r\).  Thus \(\eta\) has exact order \(n\).

Set
\[
 \xi=\operatorname{Ind}_{P}^{S_p}\eta,\qquad
 \Theta=\rho\circ\operatorname{Res}_{P}^{S_p}.
\]
In Mackey's formula, double cosets outside the normalizer have
trivial intersection with \(P\).  Their terms factor through
restriction to the trivial group, where \(\eta\) is zero since
both its lattices have coefficient rank \(p\).
The normalizer quotient is \(\mathbb F_p^\times\) and supplies
each twist once.  Hence
\[
 \Theta(\xi)
   =\sum_{a\in\mathbb F_p^\times}\rho(\alpha_a^*\eta)
   =-[I_1].
\]
Induction preserves all the mesh exact sequences, so \(n\xi=0\).
The last equality and Lemma~\ref{lem:prime-to-p-surface} show that
\(\xi\) has exact order \(n\).  Reduction and generic vanishing
commute with induction.

The detector kills every Young permutation lattice in degree \(p\).
For the partition \((p)\), its cyclic restriction is the trivial
module, whose value is zero.  In every other partition of \(p\),
each part is less than \(p\), so the Young subgroup has order
prime to \(p\).  The \(P\)-action on its cosets is free, and its
restricted permutation module is a sum of regular \(A_n[P]\)-modules.
These also have value zero.  Since \(\Theta(\xi)\ne0\), generation
fails in degree \(p\).

For an arbitrary DVR \(O\), use the same finite presentations,
coefficient splittings, and induction to define \(\xi_O\).
Its order divides \(n\), and both kernel properties were established
without completeness in Lemma~\ref{lem:prime-to-p-mesh}.
Completion carries \(\xi_O\) to the class over \(\widehat O\)
just proved to have exact order \(n\).  Thus \(\xi_O\) has
exact order \(n\) as well.  Pullback of the completed detector
provides the asserted detector over \(O\).
\end{proof}

\begin{corollary}\label{cor:prime-to-p-part-all-dvr}
Let \(O\) be any characteristic-zero DVR of residue characteristic
\(p\), and let \(N\ge2\).  Write
\[
 N=p^r d,\qquad (d,p)=1.
\]
If \(d>1\), the degree-\(p\) reduction and generic kernels for
\(A_N=O[t]/(t^N)\) contain a class of exact order \(d\), detected
by a homomorphism annihilating every Young permutation lattice.
In particular generation fails for every exponent having a prime
factor different from \(p\).
\end{corollary}

\begin{proof}
Put \(m=p^r\) and \(A_d=O[u]/(u^d)\).
The map
\[
 A_d\longrightarrow A_N,\qquad u\longmapsto t^m
\]
makes \(A_N\) finite free over \(A_d\) on
\(1,t,\ldots,t^{m-1}\).  Indeed the monomials
\(u^a t^b=t^{ma+b}\), with \(0\le a<d\) and \(0\le b<m\),
are exactly the \(O\)-basis monomials of \(A_N\).
Coefficient extension and coefficient restriction preserve
coefficient-free lattices and their coefficient-split exact sequences.
The composite of restriction after extension is \(m\) copies of the
identity functor, because the displayed coefficient basis has trivial
group action.  It therefore induces multiplication by \(m\) on \(K_0\).

Let \(\xi_d\) be the exact-order-\(d\) class of
Theorem~\ref{thm:prime-to-p-all-dvr}, and extend coefficients to
obtain \(\xi_N\).  It satisfies \(d\xi_N=0\), whereas its coefficient
restriction is \(m\xi_d\), of exact order \(d\) since \((m,d)=1\).
Thus \(\xi_N\) has exact order \(d\).
Coefficient reduction and generic extension commute with this map,
so both kernel properties persist.

Compose coefficient restriction with the detector for \(A_d\).
Restriction sends an \(A_N\)-Young lattice to \(m\) copies of the
corresponding \(A_d\)-Young lattice, and so the resulting detector
kills the Young span.  Its value on \(\xi_N\) is \(m\) times
the value on \(\xi_d\), which has exact order \(d\); this value
is nonzero.  The same proof gives an exact-order-\(\ell^{v_\ell(N)}\)
class for every prime divisor \(\ell\ne p\), by taking that maximal
prime-power divisor in place of \(d\).
\end{proof}

\begin{remark}
These arguments do not convert an order-\(d\) subgroup into a split
summand.  The remaining exponent cases not addressed here are pure
powers of the residue prime.  Combining this result with the
already established obstruction when the truncation exponent does
not divide the ramification index leaves only those pure
\(p\)-power exponents which do divide it; the dual-number case in
residue characteristic two is covered separately.
\end{remark}
